\section{LLM-Based Agent Layer}
\label{sec:llm}
The LLM layer functions as a thin wrapper around the underlying BDI agent. 
It interprets high-level textual commands from a mission controller, utilizing tools exposed by the BDI core to inject new desires, 
adjust scoring rules, or initiate team coordination. 

The mission controller can issue requests that fall into three broad categories: direct navigation goals, persistent behavioural adjustments, and team-level control. 
Each category is handled through a different mechanism in the BDI core, chosen to integrate naturally with the existing deliberation loop rather than bypassing it.

\subsection{Adapting Behavior Through Goal Injection}
The simplest requests are navigation goals: move to a tile, or deliver at a specific location.
These are translated into desires that enter the BDI intention queue and compete with existing goals through the standard scoring mechanism.
A \emph{go-to} request creates a \texttt{REACH\_TILE} desire with a configurable reward and expiry time; the reward determines how aggressively the desire competes against parcel collection.
A \emph{deliver-at} request creates a \texttt{DELIVER\_PARCEL} desire targeting a specific tile, with an optional bonus reward that biases the delivery score for that destination (Section~\ref{sec:scoring_rules}).

All those requests have a limited duration to be executed. The generated desire are made persistent for the duration of the specified TTL, so that they can be continuously re-evaluated and pursued until completion or expiry.

Both desire types use the same rate formula as their BDI counterparts (Section~\ref{sec:desires}).
A \texttt{REACH\_TILE} desire is scored as:
\begin{equation*}
  S = \frac{r_{\text{inj}} + \sum_{c \in C} r_c}{d}
\end{equation*}
where $r_{\text{inj}}$ is the configurable reward supplied by the mission controller and $\sum_{c \in C} r_c$ is the total estimated value of currently carried parcels.
This mirrors the \texttt{ReachParcel} formula, so a navigation goal with a reward equal to a typical parcel value competes on equal footing with parcel collection.
A \texttt{DELIVER\_PARCEL} desire applies the same rule-adjusted scoring as the BDI delivery desire, with an optional additive bonus applied after the rule combination, allowing the mission controller to bias delivery toward a specific tile.

\subsection{Adapting Strategy via Rule Injection}
\label{sec:scoring_rules}
The \texttt{RuleStore} maintains a persistent registry of LLM-defined scoring rules.
Each rule encodes a state condition alongside a linear effect: a multiplier $m$ and an additive offset $a$.

During each deliberation cycle, any matching rules are applied to the baseline scores from Section~\ref{sec:desires} as a linear combination. 
This scales or shifts the priorities without overwriting the foundational scoring logic.

These rules evaluate three primary state variables: the individual parcel reward, the current carry count, and the target delivery tile. 
Their effects compose directly onto the reward aggregates utilized in the desire formulas.

\paragraph{ReachParcel:}
The baseline score (Section~\ref{sec:desires}) is linearly adjusted according to the \emph{estimated stack count} post-pickup: the hypothetical carry count $|C|+1$.
This allows a rule to incentivize picking up a specific parcel only if it achieves a desired total stack size, leaving all other pickup valuations unchanged.

\paragraph{DeliverParcel:}
The delivery score is linearly modified based on the current carry count $|C|$ and, optionally, the specific target tile.
A stack rule keyed to $|C|$ dynamically alters the attractiveness of delivery depending on the agent's current payload, while a tile-specific rule can shepherd the agent toward a preferred drop-off location.
An independent bonus can also be applied directly via the \texttt{request\_putdown\_at} tool, biasing delivery toward a specific tile regardless of any active rules.

\paragraph{Explore:} 
In the absence of an active stack rule, \texttt{Explore} maintains its default priority of 0 and functions as outlined in Section~\ref{sec:desires}.

However, when a positive stack rule is triggered, \texttt{Explore} is elevated to priority 1, forcing it to compete directly against active parcel and delivery goals.
Its score is subsequently scaled by the anticipated reward of completing the stack, incentivizing the agent to actively search for missing parcels rather than settling for a premature delivery.
This behavior emerges naturally from the linear weighting system: the LLM is not required to issue an explicit exploration command, as simply registering a stack rule is sufficient to trigger the shift.

\paragraph{Traversal Penalties:}
Traversal penalties operate at the planning layer rather than the deliberation layer.
Unlike scoring rules, which modulate desire priority during deliberation, a traversal penalty injects an additive cost onto the A* edge leading onto a specific tile.
When the planner evaluates candidate paths, it adds this cost to any step landing on a penalized tile, causing the agent to prefer a detour whenever the detour length is shorter than the penalty,
without completely forbidding the option of traversing the tile if it becomes necessary.

\subsection{Coordination and Communication}
\label{sec:coordination}

\paragraph{Beacon Communication:}
To handle the critical need for up-to-date position information in coordination maneuvers, each agent continuously broadcasts its current position 
to all configured teammates as a \texttt{position\_beacon} message at a fixed interval.
This ensures that every team member maintains an up-to-date position estimate for its peers regardless of the game's sensing radius, 
which only reveals agents within a limited Manhattan distance.

\paragraph{Peer Communication:}
Simple team-level commands (navigation goals, scoring rules, traversal penalties) follow a uniform two-step pattern: the injection is applied to the issuing agent locally, then packaged as a peer-injection message and broadcast to all teammates.
Commands that carry stronger team-synchronisation requirements — rendezvous and traffic control — instead use a two-phase commit protocol described below.

\paragraph{Rendezvous and Traffic Control:}
Both \texttt{request\_rendezvous} and \texttt{request\_red\_light} employ a two-phase commit (2PC) protocol before any agent holds, ensuring the team only stops when it is collectively optimal to do so.

In Phase~1, the proposing agent broadcasts a \emph{propose} message carrying a unique round identifier (\texttt{rid}) and the command parameters.
Each receiving teammate independently evaluates whether accepting is beneficial: it computes the best hold-tile score achievable under the proposal and compares it against its current best \texttt{REACH\_PARCEL} or \texttt{DELIVER\_PARCEL} score.
For rendezvous this uses the joint-bottleneck formula (Section~\ref{sec:desires}); for traffic control it uses the standard hold-tile rate score.
The teammate replies with a \emph{vote} message containing the \texttt{rid} and its accept/reject decision.

In Phase~2, after a fixed collection window of $750 \text{ ms}$, the proposer checks for unanimity.
If every expected teammate voted \emph{accept} for the current \texttt{rid}, the proposer broadcasts a \emph{commit} message; each agent (including the proposer) then injects the \texttt{HOLD\_TILE} desires into its intention pool.
If any vote is missing or negative, the proposer broadcasts \emph{abort} and no holds are injected on any agent.

Upon commit, rendezvous holds carry a \texttt{releaseZone} that triggers automatic release once position beacons confirm all teammates are within the designated Manhattan radius — no green-light command is needed.
Traffic-control holds have no release zone and persist until their TTL expires or \texttt{request\_green\_light} is issued, which clears all active \texttt{HOLD\_TILE} intentions on both agents immediately via a fire-and-forget broadcast.
By default agents hold on the nearest odd-row tile, but the command accepts optional \texttt{axis} (\texttt{row}/\texttt{column}) and \texttt{parity} (\texttt{odd}/\texttt{even}) parameters to stagger agents across any row or column alignment, enabling flexible traffic separation strategies.

\paragraph{Coordinator:}
A centralized \texttt{Coordinator} periodically executes team assignment rounds.
During each round, a \texttt{request\_beliefs} envelope is broadcast to all teammates. After a short collection window, the coordinator constructs a geometric summary of the team's state (including agent positions, payloads, and high-value zones) and feeds it to \texttt{LLMClient.coordinate}. This triggers \texttt{assign\_goto} commands to optimally distribute the agents.
The mission controller can also trigger this process on demand via the \texttt{request\_team\_status} tool.
